% =============================================================================
% Figure 1 --- editable, \input-able TikZ source.
% Two-column band: skill library on the left, one trace on the right. ~2.8 cm tall.
%
% ---- main.tex preamble needs -------------------------------------------------
%   \usepackage{tikz}
%   \usetikzlibrary{arrows.meta}
%   \usepackage{amsmath,amssymb}
%
% ---- main.tex body -----------------------------------------------------------
%   \begin{figure}[t]
%     \centering
%     % =============================================================================
% Figure 1 --- editable, \input-able TikZ source.
% Two-column band: skill library on the left, one trace on the right. ~2.8 cm tall.
%
% ---- main.tex preamble needs -------------------------------------------------
%   \usepackage{tikz}
%   \usetikzlibrary{arrows.meta}
%   \usepackage{amsmath,amssymb}
%
% ---- main.tex body -----------------------------------------------------------
%   \begin{figure}[t]
%     \centering
%     % =============================================================================
% Figure 1 --- editable, \input-able TikZ source.
% Two-column band: skill library on the left, one trace on the right. ~2.8 cm tall.
%
% ---- main.tex preamble needs -------------------------------------------------
%   \usepackage{tikz}
%   \usetikzlibrary{arrows.meta}
%   \usepackage{amsmath,amssymb}
%
% ---- main.tex body -----------------------------------------------------------
%   \begin{figure}[t]
%     \centering
%     % =============================================================================
% Figure 1 --- editable, \input-able TikZ source.
% Two-column band: skill library on the left, one trace on the right. ~2.8 cm tall.
%
% ---- main.tex preamble needs -------------------------------------------------
%   \usepackage{tikz}
%   \usetikzlibrary{arrows.meta}
%   \usepackage{amsmath,amssymb}
%
% ---- main.tex body -----------------------------------------------------------
%   \begin{figure}[t]
%     \centering
%     \input{figures/fig1_model_inline}
%     \input{figures/fig1_caption}
%   \end{figure}
%
% Everything is namespaced \figone.../fig1... so nothing clashes with the paper.
% Cell width is DERIVED from the occurrence lists below, so the instance bars
% always line up with their occurrences: edit a sequence and the layout re-fits.
% =============================================================================

% ============================== TWEAK ZONE ===================================
% (1) The observed CPA occurrences emitted by each skill instance.
\def\figoneseqA{A,C,B,D}
\def\figoneseqB{B,C,B,C,E}
\def\figoneseqC{A,B,C,D}
\def\figoneseqD{E,B,A,C,D}

% (2) Skill colours (Okabe--Ito: blue / bluish-green / vermillion, CVD-safe).
\definecolor{fig1one}{HTML}{2E6FA7}
\definecolor{fig1two}{HTML}{2E8B6B}
\definecolor{fig1three}{HTML}{C06B2C}
\definecolor{fig1ink}{HTML}{262626}
\definecolor{fig1soft}{HTML}{8C8C8C}

% (3) Width, the left/right split, and the two row heights (cm).
\ifdefined\figonedim\else\newdimen\figonedim\fi
\ifdefined\figonefixedwidth\else\figonedim=\textwidth\fi  % \columnwidth if 2-col
\def\figonesplit{0.29}         % fraction of the width taken by the library
\def\figonegut{0.55}           % gutter between the two columns
\def\figonelabw{1.02}          % margin holding the latent/observed labels
\def\figonegp{0.15}            % gap between instances
\def\figoneyBar{-0.84}         % instance bars
\def\figoneyCell{-1.34}        % occurrence cells
% ============================ END TWEAK ZONE =================================

% ---- machinery: count the items of a comma list -----------------------------
\ifdefined\figonecnt\else\newcount\figonecnt\fi
\def\figonelistcount#1#2{%
  \global\figonecnt=0\relax
  \edef\figonelist{#2}%              % expand the \figoneseq... macro first
  \foreach \figonetmp in \figonelist{\global\advance\figonecnt by 1}%
  \xdef#1{\the\figonecnt}}

% ---- styles ------------------------------------------------------------------
\tikzset{
  fig1role/.style={inner sep=0.9pt,font=\scriptsize,text=#1},
  fig1hedge/.style={draw=fig1soft!75,line width=0.35pt,
                    -{Straight Barb[length=0.8mm,width=0.8mm]},
                    shorten >=0.6pt,shorten <=0.6pt},
  fig1name/.style={font=\footnotesize,text=#1},
  fig1ptitle/.style={anchor=west,font=\small,text=fig1ink},
  fig1gen/.style={font=\scriptsize,text=fig1ink},
  fig1rowlab/.style={font=\scriptsize,text=fig1soft,anchor=east},
  fig1note/.style={font=\scriptsize\itshape,text=fig1soft},
}

% ---- one run of observed CPA cells:  #1 start x, #2 colour, #3 letters -------
\def\figonecells#1#2#3{%
  \figonelistcount{\figoneNtmp}{#3}%
  \draw[rounded corners=1.4pt,draw=#2!28,fill=#2!8,line width=0.4pt]
    ({#1},{\figoneyCell-0.21})
    rectangle ({#1+\figoneNtmp*\figonecw},{\figoneyCell+0.21});
  \edef\figonelist{#3}%
  \foreach \figonel [count=\figonei from 0] in \figonelist{%
    \node[font=\scriptsize,text=fig1ink]
      at ({#1+(\figonei+0.5)*\figonecw},{\figoneyCell}) {\figonel};}}

% ---- one instance bar: #1 x, #2 width, #3 colour, #4 label -------------------
\def\figonebar#1#2#3#4{%
  \draw[rounded corners=1.2pt,fill=#3,draw=none]
    ({#1},{\figoneyBar-0.145}) rectangle ({#1+#2},{\figoneyBar+0.145});
  \node[font=\scriptsize,text=white] at ({#1+#2/2},{\figoneyBar}) {#4};}

\figonelistcount{\figoneNa}{\figoneseqA}
\figonelistcount{\figoneNb}{\figoneseqB}
\figonelistcount{\figoneNc}{\figoneseqC}
\figonelistcount{\figoneNd}{\figoneseqD}

\begin{tikzpicture}[x=1cm,y=1cm]

% ---- derived geometry --------------------------------------------------------
\pgfmathsetmacro{\figonew}{\the\figonedim/28.45274}
\pgfmathsetmacro{\figonelw}{\figonesplit*\figonew}          % library column
\pgfmathsetmacro{\figonelabx}{\figonelw+\figonegut}         % row-label margin
\pgfmathsetmacro{\figonerx}{\figonelabx+\figonelabw}        % trace column origin
\pgfmathsetmacro{\figonerw}{\figonew-\figonerx}
\pgfmathsetmacro{\figoneN}{\figoneNa+\figoneNb+\figoneNc+\figoneNd}
\pgfmathsetmacro{\figonecw}{(\figonerw-3*\figonegp)/\figoneN}
\pgfmathsetmacro{\figoneax}{\figonerx}
\pgfmathsetmacro{\figoneaw}{\figoneNa*\figonecw}
\pgfmathsetmacro{\figonebx}{\figoneax+\figoneaw+\figonegp}
\pgfmathsetmacro{\figonebw}{\figoneNb*\figonecw}
\pgfmathsetmacro{\figonecx}{\figonebx+\figonebw+\figonegp}
\pgfmathsetmacro{\figonecwid}{\figoneNc*\figonecw}
\pgfmathsetmacro{\figonedx}{\figonecx+\figonecwid+\figonegp}
\pgfmathsetmacro{\figonedw}{\figoneNd*\figonecw}
\pgfmathsetmacro{\figonecardi}{\figonelw/6}
\pgfmathsetmacro{\figonecardii}{\figonelw/2}
\pgfmathsetmacro{\figonecardiii}{5*\figonelw/6}

% keep the figure exactly \textwidth wide (no overfull \hbox)
\path[use as bounding box] (0,0.28) rectangle ({\figonew},-2.22);

% ============ A (left). the reusable library of skill types ==================
\node[fig1ptitle] at (0,0) {\textbf{A}\quad Reusable skill types};

% --- Skill 1: A > B, A > C, B > D, C > D   (B and C incomparable) -----------
\begin{scope}[shift={(\figonecardi,0)}]
  \node[fig1name=fig1one] at (0,-0.42) {$k{=}1$};
  \node[fig1role=fig1one] (fone) at ( 0.00,-0.84) {A};
  \node[fig1role=fig1one] (ftwo) at (-0.34,-1.28) {B};
  \node[fig1role=fig1one] (fthr) at ( 0.34,-1.28) {C};
  \node[fig1role=fig1one] (ffou) at ( 0.00,-1.72) {D};
  \draw[fig1hedge] (fone)--(ftwo); \draw[fig1hedge] (fone)--(fthr);
  \draw[fig1hedge] (ftwo)--(ffou); \draw[fig1hedge] (fthr)--(ffou);
  \node[font=\scriptsize,text=fig1soft] at (0,-1.28) {$\parallel$};
\end{scope}

% --- Skill 2: B > C > E ------------------------------------------------------
\begin{scope}[shift={(\figonecardii,0)}]
  \node[fig1name=fig1two] at (0,-0.42) {$k{=}2$};
  \node[fig1role=fig1two] (sone) at (0,-0.84) {B};
  \node[fig1role=fig1two] (stwo) at (0,-1.28) {C};
  \node[fig1role=fig1two] (sthr) at (0,-1.72) {E};
  \draw[fig1hedge] (sone)--(stwo); \draw[fig1hedge] (stwo)--(sthr);
\end{scope}

% --- Skill 3: E > {A,B,C}, C > D --------------------------------------------
\begin{scope}[shift={(\figonecardiii,0)}]
  \node[fig1name=fig1three] at (0,-0.42) {$k{=}3$};
  \node[fig1role=fig1three] (tone) at ( 0.00,-0.84) {E};
  \node[fig1role=fig1three] (ttwo) at (-0.40,-1.28) {A};
  \node[fig1role=fig1three] (tthr) at ( 0.00,-1.28) {B};
  \node[fig1role=fig1three] (tfou) at ( 0.40,-1.28) {C};
  \node[fig1role=fig1three] (tfiv) at ( 0.40,-1.72) {D};
  \draw[fig1hedge] (tone)--(ttwo); \draw[fig1hedge] (tone)--(tthr);
  \draw[fig1hedge] (tone)--(tfou); \draw[fig1hedge] (tfou)--(tfiv);
\end{scope}

\node[fig1gen,anchor=west] at (0,-2.04) {$H_k=h(U_k)$};

% ============ B (right). one trace: instances and what they emit =============
\node[fig1ptitle] at ({\figonelabx},0)
  {\textbf{B}\quad Skill instances and the CPA trace they emit};

% what is inferred vs. what is seen
\node[fig1rowlab] at ({\figonerx-0.14},{\figoneyBar})  {latent};
\node[fig1rowlab] at ({\figonerx-0.14},{\figoneyCell}) {observed};

\figonebar{\figoneax}{\figoneaw}{fig1one}{$k{=}1$}
\figonebar{\figonebx}{\figonebw}{fig1two}{$k{=}2$}
\figonebar{\figonecx}{\figonecwid}{fig1one}{$k{=}1$}
\figonebar{\figonedx}{\figonedw}{fig1three}{$k{=}3$}

\figonecells{\figoneax}{fig1one}{\figoneseqA}
\figonecells{\figonebx}{fig1two}{\figoneseqB}
\figonecells{\figonecx}{fig1one}{\figoneseqC}
\figonecells{\figonedx}{fig1three}{\figoneseqD}

% the two invocations of skill 1 differ only in the order of the incomparable pair
\draw[draw=fig1one,line width=0.6pt]
  ({\figoneax+\figonecw+0.07},{\figoneyCell-0.155})
  -- ({\figoneax+3*\figonecw-0.07},{\figoneyCell-0.155});
\draw[draw=fig1one,line width=0.6pt]
  ({\figonecx+\figonecw+0.07},{\figoneyCell-0.155})
  -- ({\figonecx+3*\figonecw-0.07},{\figoneyCell-0.155});

\node[fig1gen,anchor=west] at ({\figonerx},-1.88) {$z_n=(1,2,1,3)$};
\node[fig1note,anchor=west] at ({\figonerx+2.55},-1.88)
  {\textcolor{fig1one}{\rule[0.6ex]{7pt}{0.5pt}}\, $B\parallel C$: same type, different order};
\node[fig1gen,anchor=east] at ({\figonew},-1.88) {$\vect x^{(n)}$};

\end{tikzpicture}

%     % Caption for Figure 1.  Requires amsmath in the preamble.
\caption{\textbf{A segmental model of reusable partial-order skills.} A library
of skill types $k=1,\ldots,K$ (A), each a partial order $H_k=h(U_k)$ over CPA roles, is
reused across a trace, which is segmented into contiguous, variable-length
instances that each emit one block of observed occurrences (B). Colour marks
type: instances~1 and~3 reuse skill~1 but emit different linear extensions of it
(underlined), and instance~2 re-executes $B$ and $C$ after invalidation, so a
block need not be a linear extension of its skill at all.}
\label{fig:model}

%   \end{figure}
%
% Everything is namespaced \figone.../fig1... so nothing clashes with the paper.
% Cell width is DERIVED from the occurrence lists below, so the instance bars
% always line up with their occurrences: edit a sequence and the layout re-fits.
% =============================================================================

% ============================== TWEAK ZONE ===================================
% (1) The observed CPA occurrences emitted by each skill instance.
\def\figoneseqA{A,C,B,D}
\def\figoneseqB{B,C,B,C,E}
\def\figoneseqC{A,B,C,D}
\def\figoneseqD{E,B,A,C,D}

% (2) Skill colours (Okabe--Ito: blue / bluish-green / vermillion, CVD-safe).
\definecolor{fig1one}{HTML}{2E6FA7}
\definecolor{fig1two}{HTML}{2E8B6B}
\definecolor{fig1three}{HTML}{C06B2C}
\definecolor{fig1ink}{HTML}{262626}
\definecolor{fig1soft}{HTML}{8C8C8C}

% (3) Width, the left/right split, and the two row heights (cm).
\ifdefined\figonedim\else\newdimen\figonedim\fi
\ifdefined\figonefixedwidth\else\figonedim=\textwidth\fi  % \columnwidth if 2-col
\def\figonesplit{0.29}         % fraction of the width taken by the library
\def\figonegut{0.55}           % gutter between the two columns
\def\figonelabw{1.02}          % margin holding the latent/observed labels
\def\figonegp{0.15}            % gap between instances
\def\figoneyBar{-0.84}         % instance bars
\def\figoneyCell{-1.34}        % occurrence cells
% ============================ END TWEAK ZONE =================================

% ---- machinery: count the items of a comma list -----------------------------
\ifdefined\figonecnt\else\newcount\figonecnt\fi
\def\figonelistcount#1#2{%
  \global\figonecnt=0\relax
  \edef\figonelist{#2}%              % expand the \figoneseq... macro first
  \foreach \figonetmp in \figonelist{\global\advance\figonecnt by 1}%
  \xdef#1{\the\figonecnt}}

% ---- styles ------------------------------------------------------------------
\tikzset{
  fig1role/.style={inner sep=0.9pt,font=\scriptsize,text=#1},
  fig1hedge/.style={draw=fig1soft!75,line width=0.35pt,
                    -{Straight Barb[length=0.8mm,width=0.8mm]},
                    shorten >=0.6pt,shorten <=0.6pt},
  fig1name/.style={font=\footnotesize,text=#1},
  fig1ptitle/.style={anchor=west,font=\small,text=fig1ink},
  fig1gen/.style={font=\scriptsize,text=fig1ink},
  fig1rowlab/.style={font=\scriptsize,text=fig1soft,anchor=east},
  fig1note/.style={font=\scriptsize\itshape,text=fig1soft},
}

% ---- one run of observed CPA cells:  #1 start x, #2 colour, #3 letters -------
\def\figonecells#1#2#3{%
  \figonelistcount{\figoneNtmp}{#3}%
  \draw[rounded corners=1.4pt,draw=#2!28,fill=#2!8,line width=0.4pt]
    ({#1},{\figoneyCell-0.21})
    rectangle ({#1+\figoneNtmp*\figonecw},{\figoneyCell+0.21});
  \edef\figonelist{#3}%
  \foreach \figonel [count=\figonei from 0] in \figonelist{%
    \node[font=\scriptsize,text=fig1ink]
      at ({#1+(\figonei+0.5)*\figonecw},{\figoneyCell}) {\figonel};}}

% ---- one instance bar: #1 x, #2 width, #3 colour, #4 label -------------------
\def\figonebar#1#2#3#4{%
  \draw[rounded corners=1.2pt,fill=#3,draw=none]
    ({#1},{\figoneyBar-0.145}) rectangle ({#1+#2},{\figoneyBar+0.145});
  \node[font=\scriptsize,text=white] at ({#1+#2/2},{\figoneyBar}) {#4};}

\figonelistcount{\figoneNa}{\figoneseqA}
\figonelistcount{\figoneNb}{\figoneseqB}
\figonelistcount{\figoneNc}{\figoneseqC}
\figonelistcount{\figoneNd}{\figoneseqD}

\begin{tikzpicture}[x=1cm,y=1cm]

% ---- derived geometry --------------------------------------------------------
\pgfmathsetmacro{\figonew}{\the\figonedim/28.45274}
\pgfmathsetmacro{\figonelw}{\figonesplit*\figonew}          % library column
\pgfmathsetmacro{\figonelabx}{\figonelw+\figonegut}         % row-label margin
\pgfmathsetmacro{\figonerx}{\figonelabx+\figonelabw}        % trace column origin
\pgfmathsetmacro{\figonerw}{\figonew-\figonerx}
\pgfmathsetmacro{\figoneN}{\figoneNa+\figoneNb+\figoneNc+\figoneNd}
\pgfmathsetmacro{\figonecw}{(\figonerw-3*\figonegp)/\figoneN}
\pgfmathsetmacro{\figoneax}{\figonerx}
\pgfmathsetmacro{\figoneaw}{\figoneNa*\figonecw}
\pgfmathsetmacro{\figonebx}{\figoneax+\figoneaw+\figonegp}
\pgfmathsetmacro{\figonebw}{\figoneNb*\figonecw}
\pgfmathsetmacro{\figonecx}{\figonebx+\figonebw+\figonegp}
\pgfmathsetmacro{\figonecwid}{\figoneNc*\figonecw}
\pgfmathsetmacro{\figonedx}{\figonecx+\figonecwid+\figonegp}
\pgfmathsetmacro{\figonedw}{\figoneNd*\figonecw}
\pgfmathsetmacro{\figonecardi}{\figonelw/6}
\pgfmathsetmacro{\figonecardii}{\figonelw/2}
\pgfmathsetmacro{\figonecardiii}{5*\figonelw/6}

% keep the figure exactly \textwidth wide (no overfull \hbox)
\path[use as bounding box] (0,0.28) rectangle ({\figonew},-2.22);

% ============ A (left). the reusable library of skill types ==================
\node[fig1ptitle] at (0,0) {\textbf{A}\quad Reusable skill types};

% --- Skill 1: A > B, A > C, B > D, C > D   (B and C incomparable) -----------
\begin{scope}[shift={(\figonecardi,0)}]
  \node[fig1name=fig1one] at (0,-0.42) {$k{=}1$};
  \node[fig1role=fig1one] (fone) at ( 0.00,-0.84) {A};
  \node[fig1role=fig1one] (ftwo) at (-0.34,-1.28) {B};
  \node[fig1role=fig1one] (fthr) at ( 0.34,-1.28) {C};
  \node[fig1role=fig1one] (ffou) at ( 0.00,-1.72) {D};
  \draw[fig1hedge] (fone)--(ftwo); \draw[fig1hedge] (fone)--(fthr);
  \draw[fig1hedge] (ftwo)--(ffou); \draw[fig1hedge] (fthr)--(ffou);
  \node[font=\scriptsize,text=fig1soft] at (0,-1.28) {$\parallel$};
\end{scope}

% --- Skill 2: B > C > E ------------------------------------------------------
\begin{scope}[shift={(\figonecardii,0)}]
  \node[fig1name=fig1two] at (0,-0.42) {$k{=}2$};
  \node[fig1role=fig1two] (sone) at (0,-0.84) {B};
  \node[fig1role=fig1two] (stwo) at (0,-1.28) {C};
  \node[fig1role=fig1two] (sthr) at (0,-1.72) {E};
  \draw[fig1hedge] (sone)--(stwo); \draw[fig1hedge] (stwo)--(sthr);
\end{scope}

% --- Skill 3: E > {A,B,C}, C > D --------------------------------------------
\begin{scope}[shift={(\figonecardiii,0)}]
  \node[fig1name=fig1three] at (0,-0.42) {$k{=}3$};
  \node[fig1role=fig1three] (tone) at ( 0.00,-0.84) {E};
  \node[fig1role=fig1three] (ttwo) at (-0.40,-1.28) {A};
  \node[fig1role=fig1three] (tthr) at ( 0.00,-1.28) {B};
  \node[fig1role=fig1three] (tfou) at ( 0.40,-1.28) {C};
  \node[fig1role=fig1three] (tfiv) at ( 0.40,-1.72) {D};
  \draw[fig1hedge] (tone)--(ttwo); \draw[fig1hedge] (tone)--(tthr);
  \draw[fig1hedge] (tone)--(tfou); \draw[fig1hedge] (tfou)--(tfiv);
\end{scope}

\node[fig1gen,anchor=west] at (0,-2.04) {$H_k=h(U_k)$};

% ============ B (right). one trace: instances and what they emit =============
\node[fig1ptitle] at ({\figonelabx},0)
  {\textbf{B}\quad Skill instances and the CPA trace they emit};

% what is inferred vs. what is seen
\node[fig1rowlab] at ({\figonerx-0.14},{\figoneyBar})  {latent};
\node[fig1rowlab] at ({\figonerx-0.14},{\figoneyCell}) {observed};

\figonebar{\figoneax}{\figoneaw}{fig1one}{$k{=}1$}
\figonebar{\figonebx}{\figonebw}{fig1two}{$k{=}2$}
\figonebar{\figonecx}{\figonecwid}{fig1one}{$k{=}1$}
\figonebar{\figonedx}{\figonedw}{fig1three}{$k{=}3$}

\figonecells{\figoneax}{fig1one}{\figoneseqA}
\figonecells{\figonebx}{fig1two}{\figoneseqB}
\figonecells{\figonecx}{fig1one}{\figoneseqC}
\figonecells{\figonedx}{fig1three}{\figoneseqD}

% the two invocations of skill 1 differ only in the order of the incomparable pair
\draw[draw=fig1one,line width=0.6pt]
  ({\figoneax+\figonecw+0.07},{\figoneyCell-0.155})
  -- ({\figoneax+3*\figonecw-0.07},{\figoneyCell-0.155});
\draw[draw=fig1one,line width=0.6pt]
  ({\figonecx+\figonecw+0.07},{\figoneyCell-0.155})
  -- ({\figonecx+3*\figonecw-0.07},{\figoneyCell-0.155});

\node[fig1gen,anchor=west] at ({\figonerx},-1.88) {$z_n=(1,2,1,3)$};
\node[fig1note,anchor=west] at ({\figonerx+2.55},-1.88)
  {\textcolor{fig1one}{\rule[0.6ex]{7pt}{0.5pt}}\, $B\parallel C$: same type, different order};
\node[fig1gen,anchor=east] at ({\figonew},-1.88) {$\vect x^{(n)}$};

\end{tikzpicture}

%     % Caption for Figure 1.  Requires amsmath in the preamble.
\caption{\textbf{A segmental model of reusable partial-order skills.} A library
of skill types $k=1,\ldots,K$ (A), each a partial order $H_k=h(U_k)$ over CPA roles, is
reused across a trace, which is segmented into contiguous, variable-length
instances that each emit one block of observed occurrences (B). Colour marks
type: instances~1 and~3 reuse skill~1 but emit different linear extensions of it
(underlined), and instance~2 re-executes $B$ and $C$ after invalidation, so a
block need not be a linear extension of its skill at all.}
\label{fig:model}

%   \end{figure}
%
% Everything is namespaced \figone.../fig1... so nothing clashes with the paper.
% Cell width is DERIVED from the occurrence lists below, so the instance bars
% always line up with their occurrences: edit a sequence and the layout re-fits.
% =============================================================================

% ============================== TWEAK ZONE ===================================
% (1) The observed CPA occurrences emitted by each skill instance.
\def\figoneseqA{A,C,B,D}
\def\figoneseqB{B,C,B,C,E}
\def\figoneseqC{A,B,C,D}
\def\figoneseqD{E,B,A,C,D}

% (2) Skill colours (Okabe--Ito: blue / bluish-green / vermillion, CVD-safe).
\definecolor{fig1one}{HTML}{2E6FA7}
\definecolor{fig1two}{HTML}{2E8B6B}
\definecolor{fig1three}{HTML}{C06B2C}
\definecolor{fig1ink}{HTML}{262626}
\definecolor{fig1soft}{HTML}{8C8C8C}

% (3) Width, the left/right split, and the two row heights (cm).
\ifdefined\figonedim\else\newdimen\figonedim\fi
\ifdefined\figonefixedwidth\else\figonedim=\textwidth\fi  % \columnwidth if 2-col
\def\figonesplit{0.29}         % fraction of the width taken by the library
\def\figonegut{0.55}           % gutter between the two columns
\def\figonelabw{1.02}          % margin holding the latent/observed labels
\def\figonegp{0.15}            % gap between instances
\def\figoneyBar{-0.84}         % instance bars
\def\figoneyCell{-1.34}        % occurrence cells
% ============================ END TWEAK ZONE =================================

% ---- machinery: count the items of a comma list -----------------------------
\ifdefined\figonecnt\else\newcount\figonecnt\fi
\def\figonelistcount#1#2{%
  \global\figonecnt=0\relax
  \edef\figonelist{#2}%              % expand the \figoneseq... macro first
  \foreach \figonetmp in \figonelist{\global\advance\figonecnt by 1}%
  \xdef#1{\the\figonecnt}}

% ---- styles ------------------------------------------------------------------
\tikzset{
  fig1role/.style={inner sep=0.9pt,font=\scriptsize,text=#1},
  fig1hedge/.style={draw=fig1soft!75,line width=0.35pt,
                    -{Straight Barb[length=0.8mm,width=0.8mm]},
                    shorten >=0.6pt,shorten <=0.6pt},
  fig1name/.style={font=\footnotesize,text=#1},
  fig1ptitle/.style={anchor=west,font=\small,text=fig1ink},
  fig1gen/.style={font=\scriptsize,text=fig1ink},
  fig1rowlab/.style={font=\scriptsize,text=fig1soft,anchor=east},
  fig1note/.style={font=\scriptsize\itshape,text=fig1soft},
}

% ---- one run of observed CPA cells:  #1 start x, #2 colour, #3 letters -------
\def\figonecells#1#2#3{%
  \figonelistcount{\figoneNtmp}{#3}%
  \draw[rounded corners=1.4pt,draw=#2!28,fill=#2!8,line width=0.4pt]
    ({#1},{\figoneyCell-0.21})
    rectangle ({#1+\figoneNtmp*\figonecw},{\figoneyCell+0.21});
  \edef\figonelist{#3}%
  \foreach \figonel [count=\figonei from 0] in \figonelist{%
    \node[font=\scriptsize,text=fig1ink]
      at ({#1+(\figonei+0.5)*\figonecw},{\figoneyCell}) {\figonel};}}

% ---- one instance bar: #1 x, #2 width, #3 colour, #4 label -------------------
\def\figonebar#1#2#3#4{%
  \draw[rounded corners=1.2pt,fill=#3,draw=none]
    ({#1},{\figoneyBar-0.145}) rectangle ({#1+#2},{\figoneyBar+0.145});
  \node[font=\scriptsize,text=white] at ({#1+#2/2},{\figoneyBar}) {#4};}

\figonelistcount{\figoneNa}{\figoneseqA}
\figonelistcount{\figoneNb}{\figoneseqB}
\figonelistcount{\figoneNc}{\figoneseqC}
\figonelistcount{\figoneNd}{\figoneseqD}

\begin{tikzpicture}[x=1cm,y=1cm]

% ---- derived geometry --------------------------------------------------------
\pgfmathsetmacro{\figonew}{\the\figonedim/28.45274}
\pgfmathsetmacro{\figonelw}{\figonesplit*\figonew}          % library column
\pgfmathsetmacro{\figonelabx}{\figonelw+\figonegut}         % row-label margin
\pgfmathsetmacro{\figonerx}{\figonelabx+\figonelabw}        % trace column origin
\pgfmathsetmacro{\figonerw}{\figonew-\figonerx}
\pgfmathsetmacro{\figoneN}{\figoneNa+\figoneNb+\figoneNc+\figoneNd}
\pgfmathsetmacro{\figonecw}{(\figonerw-3*\figonegp)/\figoneN}
\pgfmathsetmacro{\figoneax}{\figonerx}
\pgfmathsetmacro{\figoneaw}{\figoneNa*\figonecw}
\pgfmathsetmacro{\figonebx}{\figoneax+\figoneaw+\figonegp}
\pgfmathsetmacro{\figonebw}{\figoneNb*\figonecw}
\pgfmathsetmacro{\figonecx}{\figonebx+\figonebw+\figonegp}
\pgfmathsetmacro{\figonecwid}{\figoneNc*\figonecw}
\pgfmathsetmacro{\figonedx}{\figonecx+\figonecwid+\figonegp}
\pgfmathsetmacro{\figonedw}{\figoneNd*\figonecw}
\pgfmathsetmacro{\figonecardi}{\figonelw/6}
\pgfmathsetmacro{\figonecardii}{\figonelw/2}
\pgfmathsetmacro{\figonecardiii}{5*\figonelw/6}

% keep the figure exactly \textwidth wide (no overfull \hbox)
\path[use as bounding box] (0,0.28) rectangle ({\figonew},-2.22);

% ============ A (left). the reusable library of skill types ==================
\node[fig1ptitle] at (0,0) {\textbf{A}\quad Reusable skill types};

% --- Skill 1: A > B, A > C, B > D, C > D   (B and C incomparable) -----------
\begin{scope}[shift={(\figonecardi,0)}]
  \node[fig1name=fig1one] at (0,-0.42) {$k{=}1$};
  \node[fig1role=fig1one] (fone) at ( 0.00,-0.84) {A};
  \node[fig1role=fig1one] (ftwo) at (-0.34,-1.28) {B};
  \node[fig1role=fig1one] (fthr) at ( 0.34,-1.28) {C};
  \node[fig1role=fig1one] (ffou) at ( 0.00,-1.72) {D};
  \draw[fig1hedge] (fone)--(ftwo); \draw[fig1hedge] (fone)--(fthr);
  \draw[fig1hedge] (ftwo)--(ffou); \draw[fig1hedge] (fthr)--(ffou);
  \node[font=\scriptsize,text=fig1soft] at (0,-1.28) {$\parallel$};
\end{scope}

% --- Skill 2: B > C > E ------------------------------------------------------
\begin{scope}[shift={(\figonecardii,0)}]
  \node[fig1name=fig1two] at (0,-0.42) {$k{=}2$};
  \node[fig1role=fig1two] (sone) at (0,-0.84) {B};
  \node[fig1role=fig1two] (stwo) at (0,-1.28) {C};
  \node[fig1role=fig1two] (sthr) at (0,-1.72) {E};
  \draw[fig1hedge] (sone)--(stwo); \draw[fig1hedge] (stwo)--(sthr);
\end{scope}

% --- Skill 3: E > {A,B,C}, C > D --------------------------------------------
\begin{scope}[shift={(\figonecardiii,0)}]
  \node[fig1name=fig1three] at (0,-0.42) {$k{=}3$};
  \node[fig1role=fig1three] (tone) at ( 0.00,-0.84) {E};
  \node[fig1role=fig1three] (ttwo) at (-0.40,-1.28) {A};
  \node[fig1role=fig1three] (tthr) at ( 0.00,-1.28) {B};
  \node[fig1role=fig1three] (tfou) at ( 0.40,-1.28) {C};
  \node[fig1role=fig1three] (tfiv) at ( 0.40,-1.72) {D};
  \draw[fig1hedge] (tone)--(ttwo); \draw[fig1hedge] (tone)--(tthr);
  \draw[fig1hedge] (tone)--(tfou); \draw[fig1hedge] (tfou)--(tfiv);
\end{scope}

\node[fig1gen,anchor=west] at (0,-2.04) {$H_k=h(U_k)$};

% ============ B (right). one trace: instances and what they emit =============
\node[fig1ptitle] at ({\figonelabx},0)
  {\textbf{B}\quad Skill instances and the CPA trace they emit};

% what is inferred vs. what is seen
\node[fig1rowlab] at ({\figonerx-0.14},{\figoneyBar})  {latent};
\node[fig1rowlab] at ({\figonerx-0.14},{\figoneyCell}) {observed};

\figonebar{\figoneax}{\figoneaw}{fig1one}{$k{=}1$}
\figonebar{\figonebx}{\figonebw}{fig1two}{$k{=}2$}
\figonebar{\figonecx}{\figonecwid}{fig1one}{$k{=}1$}
\figonebar{\figonedx}{\figonedw}{fig1three}{$k{=}3$}

\figonecells{\figoneax}{fig1one}{\figoneseqA}
\figonecells{\figonebx}{fig1two}{\figoneseqB}
\figonecells{\figonecx}{fig1one}{\figoneseqC}
\figonecells{\figonedx}{fig1three}{\figoneseqD}

% the two invocations of skill 1 differ only in the order of the incomparable pair
\draw[draw=fig1one,line width=0.6pt]
  ({\figoneax+\figonecw+0.07},{\figoneyCell-0.155})
  -- ({\figoneax+3*\figonecw-0.07},{\figoneyCell-0.155});
\draw[draw=fig1one,line width=0.6pt]
  ({\figonecx+\figonecw+0.07},{\figoneyCell-0.155})
  -- ({\figonecx+3*\figonecw-0.07},{\figoneyCell-0.155});

\node[fig1gen,anchor=west] at ({\figonerx},-1.88) {$z_n=(1,2,1,3)$};
\node[fig1note,anchor=west] at ({\figonerx+2.55},-1.88)
  {\textcolor{fig1one}{\rule[0.6ex]{7pt}{0.5pt}}\, $B\parallel C$: same type, different order};
\node[fig1gen,anchor=east] at ({\figonew},-1.88) {$\vect x^{(n)}$};

\end{tikzpicture}

%     % Caption for Figure 1.  Requires amsmath in the preamble.
\caption{\textbf{A segmental model of reusable partial-order skills.} A library
of skill types $k=1,\ldots,K$ (A), each a partial order $H_k=h(U_k)$ over CPA roles, is
reused across a trace, which is segmented into contiguous, variable-length
instances that each emit one block of observed occurrences (B). Colour marks
type: instances~1 and~3 reuse skill~1 but emit different linear extensions of it
(underlined), and instance~2 re-executes $B$ and $C$ after invalidation, so a
block need not be a linear extension of its skill at all.}
\label{fig:model}

%   \end{figure}
%
% Everything is namespaced \figone.../fig1... so nothing clashes with the paper.
% Cell width is DERIVED from the occurrence lists below, so the instance bars
% always line up with their occurrences: edit a sequence and the layout re-fits.
% =============================================================================

% ============================== TWEAK ZONE ===================================
% (1) The observed CPA occurrences emitted by each skill instance.
\def\figoneseqA{A,C,B,D}
\def\figoneseqB{B,C,B,C,E}
\def\figoneseqC{A,B,C,D}
\def\figoneseqD{E,B,A,C,D}

% (2) Skill colours (Okabe--Ito: blue / bluish-green / vermillion, CVD-safe).
\definecolor{fig1one}{HTML}{2E6FA7}
\definecolor{fig1two}{HTML}{2E8B6B}
\definecolor{fig1three}{HTML}{C06B2C}
\definecolor{fig1ink}{HTML}{262626}
\definecolor{fig1soft}{HTML}{8C8C8C}

% (3) Width, the left/right split, and the two row heights (cm).
\ifdefined\figonedim\else\newdimen\figonedim\fi
\ifdefined\figonefixedwidth\else\figonedim=\textwidth\fi  % \columnwidth if 2-col
\def\figonesplit{0.29}         % fraction of the width taken by the library
\def\figonegut{0.55}           % gutter between the two columns
\def\figonelabw{1.02}          % margin holding the latent/observed labels
\def\figonegp{0.15}            % gap between instances
\def\figoneyBar{-0.84}         % instance bars
\def\figoneyCell{-1.34}        % occurrence cells
% ============================ END TWEAK ZONE =================================

% ---- machinery: count the items of a comma list -----------------------------
\ifdefined\figonecnt\else\newcount\figonecnt\fi
\def\figonelistcount#1#2{%
  \global\figonecnt=0\relax
  \edef\figonelist{#2}%              % expand the \figoneseq... macro first
  \foreach \figonetmp in \figonelist{\global\advance\figonecnt by 1}%
  \xdef#1{\the\figonecnt}}

% ---- styles ------------------------------------------------------------------
\tikzset{
  fig1role/.style={inner sep=0.9pt,font=\scriptsize,text=#1},
  fig1hedge/.style={draw=fig1soft!75,line width=0.35pt,
                    -{Straight Barb[length=0.8mm,width=0.8mm]},
                    shorten >=0.6pt,shorten <=0.6pt},
  fig1name/.style={font=\footnotesize,text=#1},
  fig1ptitle/.style={anchor=west,font=\small,text=fig1ink},
  fig1gen/.style={font=\scriptsize,text=fig1ink},
  fig1rowlab/.style={font=\scriptsize,text=fig1soft,anchor=east},
  fig1note/.style={font=\scriptsize\itshape,text=fig1soft},
}

% ---- one run of observed CPA cells:  #1 start x, #2 colour, #3 letters -------
\def\figonecells#1#2#3{%
  \figonelistcount{\figoneNtmp}{#3}%
  \draw[rounded corners=1.4pt,draw=#2!28,fill=#2!8,line width=0.4pt]
    ({#1},{\figoneyCell-0.21})
    rectangle ({#1+\figoneNtmp*\figonecw},{\figoneyCell+0.21});
  \edef\figonelist{#3}%
  \foreach \figonel [count=\figonei from 0] in \figonelist{%
    \node[font=\scriptsize,text=fig1ink]
      at ({#1+(\figonei+0.5)*\figonecw},{\figoneyCell}) {\figonel};}}

% ---- one instance bar: #1 x, #2 width, #3 colour, #4 label -------------------
\def\figonebar#1#2#3#4{%
  \draw[rounded corners=1.2pt,fill=#3,draw=none]
    ({#1},{\figoneyBar-0.145}) rectangle ({#1+#2},{\figoneyBar+0.145});
  \node[font=\scriptsize,text=white] at ({#1+#2/2},{\figoneyBar}) {#4};}

\figonelistcount{\figoneNa}{\figoneseqA}
\figonelistcount{\figoneNb}{\figoneseqB}
\figonelistcount{\figoneNc}{\figoneseqC}
\figonelistcount{\figoneNd}{\figoneseqD}

\begin{tikzpicture}[x=1cm,y=1cm]

% ---- derived geometry --------------------------------------------------------
\pgfmathsetmacro{\figonew}{\the\figonedim/28.45274}
\pgfmathsetmacro{\figonelw}{\figonesplit*\figonew}          % library column
\pgfmathsetmacro{\figonelabx}{\figonelw+\figonegut}         % row-label margin
\pgfmathsetmacro{\figonerx}{\figonelabx+\figonelabw}        % trace column origin
\pgfmathsetmacro{\figonerw}{\figonew-\figonerx}
\pgfmathsetmacro{\figoneN}{\figoneNa+\figoneNb+\figoneNc+\figoneNd}
\pgfmathsetmacro{\figonecw}{(\figonerw-3*\figonegp)/\figoneN}
\pgfmathsetmacro{\figoneax}{\figonerx}
\pgfmathsetmacro{\figoneaw}{\figoneNa*\figonecw}
\pgfmathsetmacro{\figonebx}{\figoneax+\figoneaw+\figonegp}
\pgfmathsetmacro{\figonebw}{\figoneNb*\figonecw}
\pgfmathsetmacro{\figonecx}{\figonebx+\figonebw+\figonegp}
\pgfmathsetmacro{\figonecwid}{\figoneNc*\figonecw}
\pgfmathsetmacro{\figonedx}{\figonecx+\figonecwid+\figonegp}
\pgfmathsetmacro{\figonedw}{\figoneNd*\figonecw}
\pgfmathsetmacro{\figonecardi}{\figonelw/6}
\pgfmathsetmacro{\figonecardii}{\figonelw/2}
\pgfmathsetmacro{\figonecardiii}{5*\figonelw/6}

% keep the figure exactly \textwidth wide (no overfull \hbox)
\path[use as bounding box] (0,0.28) rectangle ({\figonew},-2.22);

% ============ A (left). the reusable library of skill types ==================
\node[fig1ptitle] at (0,0) {\textbf{A}\quad Reusable skill types};

% --- Skill 1: A > B, A > C, B > D, C > D   (B and C incomparable) -----------
\begin{scope}[shift={(\figonecardi,0)}]
  \node[fig1name=fig1one] at (0,-0.42) {$k{=}1$};
  \node[fig1role=fig1one] (fone) at ( 0.00,-0.84) {A};
  \node[fig1role=fig1one] (ftwo) at (-0.34,-1.28) {B};
  \node[fig1role=fig1one] (fthr) at ( 0.34,-1.28) {C};
  \node[fig1role=fig1one] (ffou) at ( 0.00,-1.72) {D};
  \draw[fig1hedge] (fone)--(ftwo); \draw[fig1hedge] (fone)--(fthr);
  \draw[fig1hedge] (ftwo)--(ffou); \draw[fig1hedge] (fthr)--(ffou);
  \node[font=\scriptsize,text=fig1soft] at (0,-1.28) {$\parallel$};
\end{scope}

% --- Skill 2: B > C > E ------------------------------------------------------
\begin{scope}[shift={(\figonecardii,0)}]
  \node[fig1name=fig1two] at (0,-0.42) {$k{=}2$};
  \node[fig1role=fig1two] (sone) at (0,-0.84) {B};
  \node[fig1role=fig1two] (stwo) at (0,-1.28) {C};
  \node[fig1role=fig1two] (sthr) at (0,-1.72) {E};
  \draw[fig1hedge] (sone)--(stwo); \draw[fig1hedge] (stwo)--(sthr);
\end{scope}

% --- Skill 3: E > {A,B,C}, C > D --------------------------------------------
\begin{scope}[shift={(\figonecardiii,0)}]
  \node[fig1name=fig1three] at (0,-0.42) {$k{=}3$};
  \node[fig1role=fig1three] (tone) at ( 0.00,-0.84) {E};
  \node[fig1role=fig1three] (ttwo) at (-0.40,-1.28) {A};
  \node[fig1role=fig1three] (tthr) at ( 0.00,-1.28) {B};
  \node[fig1role=fig1three] (tfou) at ( 0.40,-1.28) {C};
  \node[fig1role=fig1three] (tfiv) at ( 0.40,-1.72) {D};
  \draw[fig1hedge] (tone)--(ttwo); \draw[fig1hedge] (tone)--(tthr);
  \draw[fig1hedge] (tone)--(tfou); \draw[fig1hedge] (tfou)--(tfiv);
\end{scope}

\node[fig1gen,anchor=west] at (0,-2.04) {$H_k=h(U_k)$};

% ============ B (right). one trace: instances and what they emit =============
\node[fig1ptitle] at ({\figonelabx},0)
  {\textbf{B}\quad Skill instances and the CPA trace they emit};

% what is inferred vs. what is seen
\node[fig1rowlab] at ({\figonerx-0.14},{\figoneyBar})  {latent};
\node[fig1rowlab] at ({\figonerx-0.14},{\figoneyCell}) {observed};

\figonebar{\figoneax}{\figoneaw}{fig1one}{$k{=}1$}
\figonebar{\figonebx}{\figonebw}{fig1two}{$k{=}2$}
\figonebar{\figonecx}{\figonecwid}{fig1one}{$k{=}1$}
\figonebar{\figonedx}{\figonedw}{fig1three}{$k{=}3$}

\figonecells{\figoneax}{fig1one}{\figoneseqA}
\figonecells{\figonebx}{fig1two}{\figoneseqB}
\figonecells{\figonecx}{fig1one}{\figoneseqC}
\figonecells{\figonedx}{fig1three}{\figoneseqD}

% the two invocations of skill 1 differ only in the order of the incomparable pair
\draw[draw=fig1one,line width=0.6pt]
  ({\figoneax+\figonecw+0.07},{\figoneyCell-0.155})
  -- ({\figoneax+3*\figonecw-0.07},{\figoneyCell-0.155});
\draw[draw=fig1one,line width=0.6pt]
  ({\figonecx+\figonecw+0.07},{\figoneyCell-0.155})
  -- ({\figonecx+3*\figonecw-0.07},{\figoneyCell-0.155});

\node[fig1gen,anchor=west] at ({\figonerx},-1.88) {$z_n=(1,2,1,3)$};
\node[fig1note,anchor=west] at ({\figonerx+2.55},-1.88)
  {\textcolor{fig1one}{\rule[0.6ex]{7pt}{0.5pt}}\, $B\parallel C$: same type, different order};
\node[fig1gen,anchor=east] at ({\figonew},-1.88) {$\vect x^{(n)}$};

\end{tikzpicture}
